\textbf{Duración:} 3 horas (en aula)\\
\textbf{Enfoque:} Brief como contrato; prompt de sistema reutilizable; orquestación multi-rol con puntos HITL.

\subsubsection*{Objetivos de la sesión}
Al finalizar, la persona participante:
\begin{enumerate}
\item Completa un brief de contenido específico para material EIC.
\item Diseña un prompt de sistema con rol, restricciones, formato y criterio de parada.
\item Propone un esquema de orquestación (2--4 roles) con aprobaciones humanas.
\item Incluye instrucciones explícitas de no fabricar citas y de señalar incertidumbre.
\end{enumerate}

\subsubsection*{Agenda (resumen)}
\begin{tabularx}{\textwidth}{@{}lclX@{}}
\toprule
Bloque & Tiempo & Contenido \\
\midrule
Recap & 10 min & Integridad y temas E1 \\
Brief & 30 min & Contrato pedagógico \\
Prompt y multi-rol & 25 min & Anatomía y roles \\
Demo orquestación & 25 min & HITL y criterio de parada \\
Taller & 50 min & Brief + prompt + diagrama \\
Cierre & 20 min & Puesta en común; TI plantillas \\
\bottomrule
\end{tabularx}

\subsubsection*{Producto esperado}
Brief + prompt de sistema v0 + esquema de orquestación (E2).

\subsubsection*{Trabajo independiente relacionado}
Sesión 3: arquitectura de libro/ebook y outline. TI: práctica de plantillas (bloque B).
